\section{Discussion}\label{sec:discussion}

\subsection{Can commitment matter when ordinary predictions agree?}

The death-receptor case gives a qualified but concrete answer. The two
variants share their sustained stable-fate set. At an automatically located
shared state, even their complete sustained observable continuations agree,
yet ligand withdrawal exposes different persistent marker behavior and terminal
futures. This separates endpoint agreement from intervention-conditioned
commitment using a fixed biological interface. The analysis connects a
literature-derived mechanism to an explicit history, retained state, and
recursive-matching certificate instead of relying on relabeled terminal
mechanisms.

The scope must remain precise. Global sustained weak-trace equivalence is
refuted by a verified finite word, and the withdrawal-aware languages are
also unequal. The reverse sustained trace inclusion remains undecided.
The biological example consequently does not show that bisimulation discovers
a difference invisible to every trace-based comparison. Its value is the
systematic characterization of where matching fails and which futures remain
available under an intervention. The exact equal-trace/unequal-branching
phenomenon remains established by the accepted synthetic example; a biological
instantiation of that stronger property is still an open extension.

Nor do the markers uniquely identify a hidden commitment state. Although the
same visible prefix has different aggregated future sets in the two models,
it is compatible with \DRHistoryStates{} internal states in each. Presenting
one state-specific certificate as an unconditional prediction for every
CASP3-positive cell would erase the very nondeterminism the comparison is
intended to respect. The model's naive terminal signature after CASP3 loss
must also not be read as actual reversal of executed apoptosis. Downstream
irreversible damage, cell-type specificity and quantitative kinetics are not
established by this qualitative reconstruction.

Calzone et al. already analyzed feedback-dependent persistence under transient
ligand exposure.\cite{Calzone2010} Recovering that result is a use-case validation
of the relational workflow, not independent discovery or experimental
confirmation. The asymmetric simulation result adds an explicit account of
which continuation fails to match; the fate analysis supplies its biological
interpretation. These are complementary outputs, not interchangeable measures.

\subsection{Can changing the interface change correspondence?}

GIM addresses this separate methodological question. PN-GDDA characterizes
label-blind local Petri-net topology,\cite{Szawulak2022} whereas the executable
relations depend on labels, initial markings, and reachable choices. For the
unchanged GIM nets, weak correspondence persists at Interfaces A and B and
is lost at C, the first tested interface exposing the known terminal
distinction. This result is limited to those three declared readouts; there is
no exhaustive interface search, unique biological boundary, or globally minimal
resolution claim. The distinct terminal biology was encoded in advance, not
discovered by formal comparison.

The two cases must not be conflated. GIM changes observability while holding
the models fixed. The death-receptor study holds the interface fixed and
compares a documented mechanistic variant under sustained and withdrawal
protocols. PN-GDDA is retained for the former structural question; no uncomputed
PN-GDDA score is assigned to the death-receptor models.

\subsection{Validation, execution assumptions and limits}

The accepted mathematical development, direction conventions, and fixed-point
algorithm are unchanged. Formal regression tests establish computational
correctness under declared semantics, not biological adequacy. The new analysis
adds independent compact/tuple execution checks, a checked hidden-sink
reduction, full-model verification of the selected continuation, exact subset
counterexamples, and mCRL2 checks of the local equivalences and the sustained
counterexample word. State-space and
memory limits remain visible results, not negative relation decisions.

Unitary asynchronous updates encode possible orders, not calibrated rates.
Reachable terminal fates and marker persistence are distinguished from
inevitability under unrestricted scheduling. The selected post-withdrawal
continuations are acyclic, but no general fairness conclusion is inferred for
the much larger global graphs. Resource amendments change storage and remove
only proven irrelevant hidden readouts; they do not retune the biological
interface after inspecting relation outcomes. The prospective configuration
record is local, not an externally registered protocol or proof of absence
of analytic bias.

The HPN-DREAM/CASPOTS confrontation remains exploratory and its weak evidence
does not become stronger because the death-receptor example was added.
Likewise, the compact curated networks are not full organism-level pathways.
Future work should test other independently authored commitment models,
resolve the global sustained relations where feasible, and use state-resolved
experimental observations to challenge the predicted intervention responses.

\section{Conclusion}\label{sec:conclusion}

Endpoint-matched qualitative cell-fate models can differ in their responses to
signal withdrawal. With a fixed interface, the death-receptor feedback variants
share a reachable CASP3-positive state with exactly equivalent sustained
continuations but unequal withdrawal continuations and different terminal
futures. An automatic certificate explains the failed matching action and
preserves the simulation direction. This is a model-level reconstruction of a
published commitment phenomenon, not a new mechanism, proof of global
trace-equivalent biological branching, or wet-lab validation. GIM remains a
complementary illustration of the first correspondence loss among three
declared interfaces. Together, the cases separate what is specified through
models and readouts from what relational and reachable-future calculations
actually establish.
